\section{Coordinated changes of variables: 21, 20, and 19 integers}
\label{sec:coordinated}

The incoming \emph{Joint Sign Guards} package already supplies a
25-integer octic and equations in 21, 20, and 19 integers of degrees
20, 32, and 88. Thus reproducing the octic is a calibration step, not
a new frontier claim. The following coordinated changes of variables
give new equations in the same respective dimensions of exact degrees
\textbf{16, 18, and 40}. They use the same complete source-34 numerical
tuple and the same two determinant scales. Every transformation below
is an integer-polynomial bijection on zero sets after an integral
triangular change of coordinates.

\subsection{The fixed coefficients and the common abbreviations}

The fixed coefficients are
\[
 \beta=2^{1064921}-1,\qquad K_0,M_C,M_F,\qquad
 A=2^{1113},\quad B=2^{666104}.
\]
Here $K_0,M_C,M_F$ are the fully supplied source-34 integers;
$M_F$ is already shifted. The six numerals are bound by the packaged
coefficient manifest and reconstructed from the literal machine/table
recipe. They occupy 615883 minimal unsigned bytes. They are fixed
coefficients, never existential parameters. In particular $\beta>0$.

For the 21-variable presentation take
\[
\begin{gathered}
 J_3,C,z,f,h,r,v,w_3,s,k,y,\gamma,o,j,Y,c,I,t,p_1,p_2,p_3
 \quad\in\mathbb Z.
\end{gathered}
\]
The letter $s$ is the shifted first Pell root (\texttt{tau\_gap} in the
literal packets); it has no connection with the scale variable $v$.
Define polynomial abbreviations
\begin{gather*}
 J_0=3J_3+1,\qquad W=4(3w_3+1),\qquad q=\beta J_0+1,\qquad X=Wq,\\
 F=(K_0+W)C-z(q-1),\qquad E=XY,\qquad a=(X+1)Y,\\
 \eta=c-Yk,\qquad\zeta=k-\eta,\qquad V=fo-c,\\
 \Delta=a^2+4a+3,\qquad D=X+ac+\gamma(4a+3),\\
 R=(q-F)(q^3-q)-C(q^2-1)+(M_C+qM_F)J_0,\\
 G=AFBC-(q-AF-BC)^2,\qquad Z=\eta\zeta,\\
 T=1+p_1^2+p_2^2+p_3(p_3+1).
\end{gather*}
The alignment is typographical only: all names on the left are
polynomial abbreviations and add no unknowns.

The nine residuals of the 21-variable equation are
\begin{align*}
 R_1&=Y-(2v+1)q^3,\\
 R_2&=s^2+(2s-k)EYk-1,\\
 R_3&=D^2-\Delta c^2-1,\\
 R_4&=k-1-hE-R,\\
 R_5&=V-jc+k-1-hE,\\
 R_6&=I-(2r+1)c^2,\\
 R_7&=I^2-\Delta(f^2-1),\\
 R_8&=y^2-1-I^2(y^2-V^2),\\
 R_9&=GZ-(t^2-G-Z+1)^2-T.
\end{align*}
Set $Q_{21}=\sum_{i=1}^9R_i^2$. The displayed degrees are,
in this order,
\[
 4,6,8,5,4,3,8,6,8.
\]
In the highest-degree part of $Z=(c-Yk)(k-c+Yk)$, the monomial
$-Y^2k^2$ occurs with coefficient $-1$. The coefficient of $Y^4k^4$
in $R_9$ is consequently $-1$, and no other residual contributes
$Y^8k^8$ to its square. Thus
\[
 [Y^8k^8]Q_{21}=1,
\]
and the exact degree is sixteen, including after all fixed ports are
specialized to their actual values.

\subsection{Why these coordinates are effective}

The literal incoming 25-integer packet uses
$k_0=1-V+jc+hE$ in three arithmetic rows, while retaining
$\eta+\zeta$ in its quotient and positivity rows. To obtain the convenient
uniform presentation used here, put
$\varepsilon=V-jc+\eta+\zeta-1-hE$. Add $(E-a)\varepsilon$ to its main-ordinate row,
add $-E(c-\eta)\varepsilon$ to its first-norm row, and add $\varepsilon$ to its packed-rank
row. Keep $\varepsilon$ itself. These are reversible elementary row operations,
so the integer zero set is unchanged. The exact checker compares all
thirteen resulting rows with the independently reconstructed packet.

In this presentation of the source 25-integer octic use the integral triangular changes
\[
 C=q-F-\alpha,\qquad Y=a-E,\qquad k=\eta+\zeta,
 \qquad s=\tau-E(c-\eta).
\]
Their inverses are
\[
 \alpha=q-F-C,\quad a=E+Y,\quad\zeta=k-\eta,
 \quad\tau=s+E(c-\eta).
\]
The source rows then determine four coordinates monically:
\[
 F=(K_0+W)C-z(q-1),\qquad E=XY,\qquad
 \eta=c-Yk,\qquad V=fo-c.
\]
Delete the four defining rows and substitute these four integer
polynomials everywhere. This reduces 25 coordinates and 13 residuals
to 21 coordinates and 9 residuals. The first norm is controlled by
the exact identity, with $u=E(c-\eta)=EYk$,
\[
 (s+u)^2-(k+u)^2+(k+u)k-1
   =s^2+(2s-k)u-1.
\]
The cancellation prevents a needless squaring of the degree-five
expression $EYk$.

Retaining both $Y$ and $k$ is essential to the bound. It keeps
$\eta,\zeta$ quadratic, even after $E$ and $F$ have been removed.
The determinant $G$ and product $Z$ are both quartic, so the combined
positivity residual remains of degree eight. Expanding the quotient
coordinates too early was responsible for the larger degree of the
inspected incoming projections.

For complete rigor, map an old zero forward by the four triangular
definitions and then forget $F,E,\eta,V$. The four deleted source
equations imply precisely the restoration polynomials above. Conversely,
restore these coordinates at any new zero and invert the triangular
changes; every source residual vanishes, including the four deleted
defining equations. The two maps are inverse on zero sets. There is
no division, denominator condition, or sign assumption in this argument.
Therefore the inherited equivalence with finite ZFC refutability transfers
to $Q_{21}$.

\subsection{Twenty integers and degree eighteen}

In $Q_{21}$ make the integral change $y=fo+b$, where $b$ is unrestricted,
and eliminate $I=(2r+1)c^2$ using $R_6$. All other coordinates remain.
The key cancellation is
\[
 (fo+b)^2-(fo-c)^2=(c+b)(2fo+b-c).
\]
The resulting auxiliary residual is
\[
 (fo+b)^2-1-(2r+1)^2c^4(c+b)(2fo+b-c),
\]
of degree nine. The seven other residuals have degrees at most eight.
The final equation $Q_{20}=0$ is their sum of squares, has twenty integer
coordinates, and has exact degree eighteen because
\[
 [r^4c^{10}f^2o^2]Q_{20}=64.
\]
Only the auxiliary row contributes this monomial. The variable change and
monic elimination are reversible over the integers. Hence $Q_{20}$ has a
zero exactly when $Q_{21}$ does.

\subsection{Nineteen integers and degree forty}

Finally substitute
\[
 Y=(2v+1)\bigl(\beta(3J_3+1)+1\bigr)^3
\]
in $Q_{20}$ and remove $R_1$. This leaves the nineteen coordinates
\[
 J_3,C,z,f,h,r,v,w_3,s,k,b,\gamma,o,j,c,t,p_1,p_2,p_3.
\]
Use exactly the common abbreviations above, the additional definitions
$I=(2r+1)c^2$ and $y=fo+b$, and the seven residuals
$R_2,R_3,R_4,R_5,R_7,R_8,R_9$. These are completely specified ordinary
polynomials, with respective degrees
\[
 12,14,7,7,14,9,20.
\]
Their squared sum is $Q_{19}$. Its exact degree is forty, since
\[
 [v^8J_3^{24}k^8]Q_{19}=2^8\,3^{24}\,\beta^{24}>0.
\]
This coefficient is the image of the $Y^8k^8$ certificate. Every other
residual has degree at most fourteen before squaring, so it cannot alter
the degree-forty coefficient. The monic restoration of $Y$ again gives a
bijection of integer zero sets.

\subsection{Proof boundary and measured comparison}

Each new system therefore has a zero exactly when the existing
25-integer octic has a zero. Combining with its named native/compiler,
signed-Pell, determinant and K215/refutation calibration gives, by ordinary
mathematics formalizable in ZFC,
\[
 \mathrm{ZFC}\ \vdash\quad
 \neg\exists\boldsymbol{x}\in\mathbb Z^n\,
       Q_n(\boldsymbol{x})=0
 \quad\longleftrightarrow\quad\operatorname{Con}(\mathrm{ZFC}),
 \qquad n\in\{19,20,21\}.
\]
No absence of zeros or consistency assertion is proved here. No new
end-to-end proof-assistant certificate or internally encoded derivation
in the repository's object-language ZFC calculus is claimed.

The comparison includes all six incoming archives visible at the pinned
repository snapshot, not just numbered proposals. The changes in
$(n,\deg)$ are
\[
 (21,20)\longrightarrow(21,16),\qquad
 (20,32)\longrightarrow(20,18),\qquad
 (19,88)\longrightarrow(19,40).
\]
These are degree reductions of 20\%,
43.75\%, and approximately 54.55\% at unchanged dimensions. Other
measures may increase; fixed coefficient payload size is unchanged.
Existing ten-integer and nine-positive high-degree equations remain
different tradeoffs. No worldwide priority or global minimum is asserted.

Exact finite checks verify every defining-row identity, both integral
root translations, all compact-display residuals, all variable counts and
liveness, the stated residual degree bounds, and the three nonzero
coefficient certificates. They also reconstruct the complete numerical
tuple. These computations audit the literal algebra; the zero-set maps
above supply the universal proof.
